%==============================================================================
\subsection{MessagePack}
%==============================================================================

\subsubsection{Example}

MessagePack data is been encode to binary data and decoded to json data.

This is how a HTTP packet would look like with MessagePack data inside:

\begin{lstlisting}[style=hexdump]
0000   00 00 23 01 04 00 00 00 05 83 04 94 60 00 00 00   ..#.........`...
0010   06 07 59 98 b5 05 b1 89 8a 98 88 7b 6a 4d 9e bf   ..Y........{jM..
0020   87 c3 c2 c1 c0 bf be 0f 0d 02 37 36 00 00 4c 00   ..........76..L.
0030   01 00 00 00 05 92 84 a7 74 69 74 6c 65 43 64 a6   ........titleCd.
0040   30 32 35 33 34 38 a6 75 73 65 72 49 64 b2 36 38   025348.userId.68
0050   35 31 32 39 32 36 30 35 32 32 31 39 32 35 34 39   5129260522192549
0060   a7 73 65 73 73 69 6f 6e ad 36 61 31 30 61 64 62   .session.6a10adb
0070   64 36 35 61 65 63 a8 70 6c 61 74 66 6f 72 6d 03   d65aec.platform.
0080   90                                                .
\end{lstlisting}

The binairy MessagePack data starts at 0x92.

You can copy this data as a hex stream in wireshark as example:

\begin{lstlisting}[style=hexstream]
9284a77469746c654364a6303235333438a6757365724964b236383531323932
3630353232313932353439a773657373696f6ead366131306164626436356165
63a8706c6174666f726d0390
\end{lstlisting}

And you can decode it to:

\begin{lstlisting}
[
  {
    "titleCd": "025348",
    "userId": "685129260522192549",
    "session": "6a10adbd65aec",
    "platform": 3
  },
  []
]
\end{lstlisting}

\begin{tcolorbox}[colback=green!10,colframe=blue!80!black,title=Info]
You can use this online tool to decode MessagePack data: https://ref45638.github.io/msgpack-converter/
\end{tcolorbox}

\newpage

\subsubsection{Request}

A recurring request object was observed throughout the API\@.

The general form is:

\begin{lstlisting}
[
    {
        "titleCd": "...",
        "userId": "...",
        "session": "...",
        "platform": 3,
        "version": "09.01"
    },
    <endpoint-specific arguments>
]
\end{lstlisting}

Not every request contains the \texttt{version} member.  The earlier
\texttt{/000000/} system and user APIs use a smaller common object.

For example:

\begin{lstlisting}
[
    {
        "titleCd": "025348",
        "userId": "685129260522192549",
        "session": "6a10adb...",
        "platform": 3
    },
    []
]
\end{lstlisting}

For the game-specific \texttt{/025348/} APIs, the observed common object is:

\begin{lstlisting}
[
    {
        "titleCd": "025348",
        "userId": "685129260522192549",
        "session": "6a10ad...",
        "platform": 3,
        "version": "09.01"
    },
    []
]
\end{lstlisting}

The second element of the outer array contains endpoint-specific arguments.
Several endpoints use an empty array, indicating that no additional
endpoint arguments were necessary in those requests.

\newpage

\subsubsection{Response}

The observed responses share a common outer structure:

\begin{lstlisting}
[
    {
        "result": 0,
        "date": "...",
        ...
    },
    <endpoint-specific data>
]
\end{lstlisting}

\newpage